% =====================================================================
% JSAC2 v1 — first true IEEE 2-column draft
% Reflective Intelligent Human Bodies: a closed-loop body digital twin
% for application-aware mmWave capacity and live SAR readings.
%
% Author: Robin Wydaeghe
% Target: IEEE JSAC SI on Digital Twins for Wireless Networks (Q4 2026)
% Spine: ../JSAC_SPINE_PLAN_v2.md
% =====================================================================

\documentclass[journal,twocolumn,10pt]{IEEEtran}

\usepackage[utf8]{inputenc}
\usepackage[T1]{fontenc}
\usepackage{lmodern}
\usepackage{microtype}

\usepackage{amsmath,amssymb,amsthm,mathtools}
\usepackage{bm}
\usepackage{bbm}

\usepackage{graphicx}
\graphicspath{{figures/}{../code/outputs/}{../code/outputs/munich/}{../../theory/figures/}{../../JSAC/code/experiments/plaza_run/figures/}}
\usepackage{booktabs}
\usepackage{array}
\usepackage{multirow}
\usepackage{caption}
\usepackage{subcaption}

\usepackage{cite}
\usepackage{algorithm}
\usepackage{algpseudocode}
\usepackage{xcolor}
\usepackage[colorlinks=true,linkcolor=blue,citecolor=blue,urlcolor=blue]{hyperref}
\usepackage[capitalize,nameinlink]{cleveref}

% Placeholder for limbo content (red box, easy to scan)
\newcommand{\limbo}[1]{{\color{red!70!black}\fbox{\parbox{0.92\columnwidth}{\textbf{TODO:} #1}}}}

\theoremstyle{plain}
\newtheorem{theorem}{Theorem}
\newtheorem{proposition}[theorem]{Proposition}
\newtheorem{corollary}[theorem]{Corollary}
\newtheorem{lemma}[theorem]{Lemma}
\theoremstyle{definition}
\newtheorem{definition}[theorem]{Definition}
\theoremstyle{remark}
\newtheorem{remark}[theorem]{Remark}

\crefname{proposition}{Proposition}{Propositions}
\Crefname{proposition}{Proposition}{Propositions}
\crefname{lemma}{Lemma}{Lemmas}
\Crefname{lemma}{Lemma}{Lemmas}
\crefname{theorem}{Theorem}{Theorems}
\Crefname{theorem}{Theorem}{Theorems}
\crefname{corollary}{Corollary}{Corollaries}
\Crefname{corollary}{Corollary}{Corollaries}
\crefname{remark}{Remark}{Remarks}
\Crefname{remark}{Remark}{Remarks}
\crefname{definition}{Definition}{Definitions}
\Crefname{definition}{Definition}{Definitions}

\newcommand{\khat}{\hat{\bm{k}}}
\newcommand{\nhat}{\hat{\bm{n}}}
\newcommand{\rhat}{\hat{\bm{r}}}
\newcommand{\etahat}{\hat{\bm{\eta}}}
\newcommand{\rr}{\mathbf{r}}
\newcommand{\xx}{\mathbf{x}}
\newcommand{\hh}{\bm{h}}
\newcommand{\EE}{\mathbf{E}}
\newcommand{\HH}{\mathbf{H}}
\newcommand{\Sinc}{S_{\mathrm{inc}}}
\newcommand{\Sab}{S_{\mathrm{ab}}}
\newcommand{\Pabs}{P_{\mathrm{ab}}}
\newcommand{\Aab}{A_{\mathrm{ab}}}
\newcommand{\Qmat}{\mathbf{Q}}
\newcommand{\Qabs}{\mathbf{Q}_{\mathrm{ab}}}
\newcommand{\Mmat}{\mathbf{M}}
\newcommand{\Jmat}{\mathbf{J}}
\newcommand{\Phimat}{\mathbf{\Phi}}
\newcommand{\Fn}{\mathbf{F}_n}
\newcommand{\psin}{\bm{\psi}_n}
\newcommand{\thetab}{\bm{\theta}}
\newcommand{\betab}{\bm{\beta}}
\newcommand{\gammab}{\bm{\gamma}}
\newcommand{\hBS}{\bm{h}_{\mathrm{BS}}}
\newcommand{\hbody}{\bm{h}_{\mathrm{b}}}
\newcommand{\Tepskin}{T_{0}}
\newcommand{\diff}{\mathrm{d}}
\newcommand{\Complex}{\mathbb{C}}
\newcommand{\Reals}{\mathbb{R}}
\newcommand{\ntilde}{\tilde{n}}
% (was: ReLU operator -- replaced by [\cdot]_+ notation throughout)
\DeclareMathOperator{\RE}{Re}
\DeclareMathOperator{\IM}{Im}
\DeclareMathOperator{\tr}{tr}
\DeclareMathOperator*{\argmin}{arg\,min}
\DeclareMathOperator*{\argmax}{arg\,max}

\title{Reflective Intelligent Human Bodies:\\
A Closed-Loop Body Digital Twin for mmWave Capacity\\
and Live Exposure Readings}

\author{Robin~Wydaeghe%
  \thanks{R. Wydaeghe is with Ghent University and imec, Ghent, Belgium
  (e-mail: robin.wydaeghe@ugent.be).}}

\markboth{IEEE Journal on Selected Areas in Communications,
  Vol.~XX, No.~X, Special Issue on Digital Twins for Wireless Networks, Month~2026}%
{Wydaeghe: Reflective intelligent human bodies}

\begin{document}

\maketitle

\begin{abstract}
Two persistent frictions slow millimetre-wave deployment, and both
locate at the same point in the propagation graph: bodies in the
line-of-sight cost 25 to 40~dB of received power, and public concern
about non-ionising exposure carries the second cost as a deployment
friction that compliance numbers alone have not dissolved. This paper
proposes that an on-device digital twin of the user's body, calibrated
in real time by uplink pilots and embedded inside the radio control
loop, addresses both. The twin is a Skinned Multi-Person Linear (SMPL-X)
parametric mesh whose pose vector $\thetab \in \Reals^{66}$ is the only
state the user can actuate. Two semantic outputs are produced from the
same per-triangle Kirchhoff--physical-optics machinery: (i)~a live
absorbed-power-density reading that, displayed locally and aggregated
opt-in across users, produces the per-individual exposure time series
that prospective RF-EMF cohort studies currently approximate via
self-report; and (ii)~a pose-differentiable cascaded mmWave channel
$\hh(\thetab) = \hBS + \hbody(\thetab)$ against which the base
station issues low-bandwidth pose suggestions to the user via the
downlink control channel. The user actuates, the IMU and uplink pilots
resolve the new state, the body-side calibration coefficient
$\gammab \in \Complex^{K}$ absorbs the residual via a $K$-mode
truncated SVD of the body-to-phone Kirchhoff operator with
$K = K_{99}$ chosen adaptively from the operator spectrum, and the
loop closes. The central new primitive is a UE-anchored Kirchhoff
render that constructs $\hbody(\thetab)$ in $\mathcal{O}(M\,T_{\mathrm{vis}})$
operations per pose by integrating per-triangle reflected fields against
the phone position rather than tracing forward radiosity rays into the
half-space, a complexity-class win when the receiver count is small.
An empirical study over $18$ Sionna~RT Munich scenes ($9$ LOS,
$9$ NLOS) reports a naturalness-bounded pose move worth $0.6$~to
$8.6\,$dB of cascaded-SINR gain (mean $1.7\,$dB on LOS, $4.1\,$dB
on NLOS), with up to $8.6\,$dB body-RIS rescue on a deep-shadow
NLOS scene (scene loss $> 100\,$dB); a $K$-iteration projected
gradient ascent on the VPoser latent (one autodiff backward per
step through VPoser, SMPL-X and the Kirchhoff render) reaches
within $1\,$dB of the brute-force ball optimum on $8/9$ LOS and
$8/9$ NLOS scenes at $K = 20$;
the gradient mass concentrates on six torso joints out of 22; and the
peak-local absorbed power density at the canonical FR2 operating point
sits at 0.58\,\% of the ICNIRP general-public limit, two orders of
magnitude below the regulatory ceiling.
\end{abstract}

\begin{IEEEkeywords}
Digital twin, twin-in-the-loop optimisation, closed-loop control,
application-aware networks, millimetre wave, MIMO, parametric body
model, SMPL-X, physical optics, Kirchhoff scattering, human-in-the-loop
control, ICNIRP, absorbed power density, electromagnetic-field
exposure, RF-EMF epidemiology.
\end{IEEEkeywords}

\IEEEpeerreviewmaketitle


% =====================================================================
\section{Introduction}\label{sec:intro}
% =====================================================================

\IEEEPARstart{T}{he} millimetre-wave band (FR2, $24$ to $52.6\,$GHz)
offers the wide bandwidths that 5G and 6G need to deliver multi-Gbps
peak rates~\cite{general5G,TS38214}. Multi-element antenna arrays at
the base station, increasingly at densities of tens to hundreds of
elements per panel, are what makes that promise reach the user in
practice~\cite{HeathLozanoFA2018}.

mmWave deployment nevertheless remains challenging. First, a single
human body sitting in the line-of-sight of an FR2 channel costs $25$
to $40\,$dB of received power~\cite{MacCartney2017,Maccartney2017blockage};
the array gain that the panel buys back is small once a pedestrian
stands between the antenna and the user. Second, public concern over
non-ionising electromagnetic exposure has translated into precautionary
caps that effectively block in-region mmWave: the Brussels Capital
Region with a strict cumulative outdoor field
cap~\cite{BrusselsArrete2007,BrusselsArrete2024}, the Italian 2024
decree with $24$-hour median attention values~\cite{ItalyDL}, and
similar Geneva-canton limits~\cite{SwissNISV2000}. Both routes pass through the
same object: the human body in the cell.

The exposure-aware beamforming literature already takes the per-body
absorbed power as a Hermitian quadratic form
$\Pabs = \xx^{H}\Qabs\xx$ in the precoder $\xx$, with the SAR matrix
$\Qabs$ a body-and-pose-dependent
state~\cite{Hochwald2014,Ying2015,Ying2017}. Reconfigurable intelligent
surfaces, in turn, model the cascaded channel that a passive scatterer
inserts between the BS and the UE~\cite{DiRenzo2020,Wu2019}.
\emph{The human body is a passive lossy reflector mediated by its pose
and actuated by intelligent instructions on a smartphone. We aptly
name this a Reflective Intelligent Human Body (RIHB).} Where a panel
RIS tunes its surface electronically, the RIHB tunes by pose.

The user is in the loop. They hold the phone, they cause the
blockage, and they are the only available actuator on the body
geometry. Yet today the user has no idea they are blocking the
signal and no signal from the network on how their pose could
improve the link. Three gaps therefore remain open.
\emph{(i)~No current body digital twin gives the public, regulators
or epidemiologists both an accurate and an instantaneous reading of
the absorbed power on the body surface.}
\emph{(ii)~The body-mediated channel from the BS to the UE is
treated, when at all, as a forward-scattering problem; the
phone-on-body geometry is a single-receiver problem and asks for
a UE-anchored render instead.}
\emph{(iii)~A user-in-the-loop closed loop --- BS gradient on a
pose-differentiable rate, sub-degree pose cue to the user, IMU
+ uplink-pilot resolution of the new state, repeat --- has not
been put together.}

This work fills the three gaps with an on-device digital twin of
the user's body that operates as a wireless-native intelligent
agent inside the per-slot control loop. The twin is a parametric
Skinned Multi-Person Linear body (SMPL-X)~\cite{SMPLX2019}; the
phone, the BS, and a Sionna-RT~\cite{Hoydis2023Sionna} scene
between them carry the propagation. From the same per-triangle
Fresnel-physical-optics machinery the twin produces two outputs:
a live absorbed-power-density reading on the body surface, and a
pose-differentiable cascaded channel against which the BS issues a
small pose suggestion at every tick. The user actuates, the IMU
and the uplink pilots resolve the new state, the body-side
calibration coefficient $\gammab$ absorbs the residual via a
$K$-mode SVD of the body-to-phone Kirchhoff operator, and the loop
proceeds. The user moves from a bad pose into a high-rate window
and exploits the bandwidth that FR2 was always meant to deliver,
while the phone shows them, for the first time, what the dose on
their body is at that moment.

To the authors' knowledge, this work is, for the first time:
\begin{enumerate}
\item a body digital twin that operates inside a per-slot
bidirectional control loop on a real RT-traced mmWave scene
(\cref{sec:arch});
\item a UE-anchored Kirchhoff render of the body-mediated channel
that is fast enough on consumer-grade compute to run at
$\sim\!10\,$Hz, where the standard FDTD reference takes
minutes-to-hours per pose on the standard IT'IS Virtual
Population FDTD reference~\cite{Christ2010VF} (\cref{sec:kirchhoff});
\item an end-to-end pose-differentiable rate, with a
human-in-the-loop control law that turns the BS gradient into a
sub-degree pose cue (\cref{sec:control});
\item the first instantaneous on-device read-out of the
absorbed-power-density field across the whole mmWave band,
$\sim\!10^{4}\times$ faster than FDTD at comparable accuracy on
biological tissue (\cref{sec:saridx}); and
\item an empirical convergence study on the canonical Sionna RT
Munich scene with $30$ Rx positions ($16$ LOS, $14$ NLOS),
reporting existence and reachability of naturalness-bounded pose
gains and a four-class failure-mode taxonomy
(\cref{sec:numerical}).
\end{enumerate}

\begin{table}[!t]
\caption{Two semantic outputs from one piece of per-triangle
Kirchhoff machinery.}
\label{tab:two-products}
\centering\small
\begin{tabular}{@{}p{0.30\columnwidth}p{0.27\columnwidth}p{0.30\columnwidth}@{}}
\toprule
& \textbf{Live exposure reading} & \textbf{Pose-driven capacity} \\
\midrule
Computes & $\Sab(\rr)$ on body surface & $\nabla_{\thetab} R(\thetab,\xx)$ \\
Output & dose number, displayed locally & pose move, suggested via DL \\
Where it lives & on the body surface & at the phone antenna \\
Privacy delta & reading local; opt-in aggregation & gradient local; only $\thetab^{*}$ accept/reject leaves the device \\
Loop role & informational & rate maximisation under comfort budget \\
\bottomrule
\end{tabular}
\end{table}


% =====================================================================
\section{System model and DTN architecture}\label{sec:arch}
% =====================================================================

\subsection{Geometry and channel}

The base station is a rectangular array of $M$ elements with
half-wavelength spacing, transmitting at carrier $f_c$ in the FR2
band ($f_c = 28$~GHz throughout, except where indicated). The user
is a single served terminal at world position $\rr_p$ holding a
phone with a single receive antenna, isotropic in pattern for
simplicity.\footnote{Polarisation- and pattern-aware extensions are
straightforward via the projection $\bm{g}_{\mathrm{UE}}^{H}$
introduced in \cref{sec:kirchhoff}; we suppress them in the main
derivation.} The cascaded narrowband channel at the phone is
\begin{equation}\label{eq:cascaded}
\hh(\thetab,\rr_p) = \hBS(\rr_p) + \hbody(\thetab,\rr_p),
\end{equation}
where $\hBS$ collects every BS-to-phone propagation path that does
not interact with the body (direct LOS when present, plus all
scene-side multipath via building reflections and edge diffraction)
and $\hbody$ is the body-mediated coefficient computed by the
on-device twin. The split is exact by Maxwell linearity and is the
operating decomposition used throughout.
The BS forms a precoder $\xx \in \Complex^{M}$ with $\|\xx\|^2 \leq
P_{\mathrm{tx}}$ and the received signal at the phone is $y = \hh^{H}
\xx s + n$ with $n \sim \mathcal{CN}(0, N_0)$. The achievable rate per Hz is
$R = \log_2(1 + \mathrm{SINR})$ with SINR defined in
\cref{sec:control}; we report rates in bps/Hz throughout the
numerics, leaving any modulation-coding-scheme cap implicit at the
implementation layer.

\subsection{Body twin and pose telemetry}

The user's body is represented as an SMPL-X parametric
mesh~\cite{SMPLX2019} with 22 body joints, 10\,475 vertices, and
20\,908 triangles. The pose state is the axis-angle vector $\thetab
\in \Reals^{22 \times 3}$, with body-shape parameters $\bm{\beta}$
(distinct from the BS-side calibration coefficient $\betab$ of
\cref{sec:closed-loop}) treated as fixed user attributes calibrated
once at enrolment. The mesh map $\thetab \mapsto
\Sigma(\thetab)$ is differentiable through linear blend skinning,
which is what makes the gradient $\nabla_{\thetab} R$ of
\cref{sec:control} a usable object.

Pose telemetry is supplied by an in-pocket consumer-grade IMU
running a Madgwick-style attitude
filter~\cite{Madgwick2011}, with steady-state per-joint attitude
RMS error $\sigma_{\mathrm{joint}} \approx 4^{\circ}$ for the
canonical operating point. The IMU error model and the
$\sigma_{\mathrm{joint}}$ sweep are deferred to \cref{sec:closed-loop}
and the supplementary information; here it suffices that pose enters the
twin via the IMU stream and is corrected over the air by the
closed-loop calibration of \cref{sec:closed-loop}.

\subsection{The closed loop}

The closed-loop architecture is shown in \cref{fig:dtn-arch}.

\begin{figure*}[!t]
\centering
\includegraphics[width=0.92\textwidth]{figures/fig_dtn_architecture.pdf}
\caption{Network-centric DTN architecture. The base station maintains
the per-element path dictionary $\Jmat$ and the precoder $\xx$
solver. The on-device body twin holds the SMPL-X mesh
$\Sigma(\thetab)$, runs the UE-anchored Kirchhoff render to
obtain $\hbody(\thetab,\bm{r}_p)$, fits the body-side calibration
coefficient $\gammab$ over a $K_{99}$-mode SVD basis from
uplink-pilot residuals, and produces both the rate gradient
$\nabla_{\thetab} R$ and the absorbed-power-density field
$\Sab(\bm{r})$. The user is the actuator of $\thetab$ and the
substrate of the IMU. Solid arrows are forward (downlink); dashed
arrows are return (uplink). Two semantic outputs land below their
producers: a network-side rate adaptation and a user-side live
exposure reading.}
\label{fig:dtn-arch}
\end{figure*}

In words, at each twin tick:
\begin{enumerate}
\item The base station refreshes the path dictionary $\Jmat$ at
$\sim\!1$~Hz; per-path amplitude calibration $\betab$ runs over uplink
pilots at the same cadence.
\item The on-device twin reads IMU pose $\hat\thetab$, runs the SMPL-X
forward map to obtain the mesh, evaluates the UE-anchored Kirchhoff
render of \cref{sec:kirchhoff} to obtain $\hbody$, and combines with
$\hBS$ to produce a cascaded-channel estimate.
\item The BS observes the served user's actual channel via uplink
pilots, computes the residual against the twin's prediction, and fits
the body-side coefficient $\gammab$ over a $K$-mode SVD basis of the
body-to-phone Kirchhoff operator (\cref{sec:closed-loop}).
\item The BS computes a pose-direction suggestion $\nabla_{\thetab}
R(\thetab) / \|\nabla_{\thetab} R\|$ truncated to a comfortable
magnitude, and emits a low-bandwidth packet over the downlink
control channel ($\sim 6 M r$ float scalars per pose suggestion at
low rank $r$; see \cref{sec:bandwidth}).
\item The user, optionally, actuates a small pose move toward the
suggestion. The loop closes at the next twin tick.
\end{enumerate}
The on-device twin separately produces the absorbed-power-density
reading $\Sab(\rr;\thetab,\xx)$ for display to the user and (opt-in)
upload to a long-running RF-EMF cohort study; this branch is
informational and does not enter the rate-optimisation loop.

\subsection{Positioning against the digital-twin literature}

Most digital-twin work in wireless networks twins the network
infrastructure itself, an industrial process, or an autonomous
agent~\cite{Chen2024DTN, Saad2025DTN, Hashash2025DTN}. The
contribution here is structurally different: the twin is the human
body in the cell, not the network or the application. The
semantic-feedback channel (the live exposure reading as
application output) and the differentiable-control channel (the
pose gradient as agent action) are first-class outputs of the same
physical model. Both directions of the loop close on the same
machinery, the per-triangle Kirchhoff render of
\cref{sec:kirchhoff}.


% =====================================================================
\section{RIHB physics: per-triangle Fresnel and the body PO scatterer}\label{sec:physics}
% =====================================================================

This section assembles the per-triangle scattering primitive that
both products of \cref{tab:two-products} consume. We give only the
narrow slice needed for what follows: a single-bounce
Fresnel-physical-optics treatment on a lossy half-space at mmWave.
Polarisation-aware Stokes-vector dosimetry, curvature corrections,
diffraction smoothing, and inter-body radiosity extensions live in
the supplementary information.

\subsection{Fresnel coefficients at a lossy half-space}

\begin{figure}[!t]
\centering
\includegraphics[width=0.95\columnwidth]{incidence_plane_fig.pdf}
\caption{Per-triangle incidence-plane geometry on the body surface
$\Sigma$. The TE basis $\hat{e}_{s,n}$ (green) is orthogonal to the
incidence plane; the incident TM basis $\hat{e}_{p,n}$ (orange) lies
in the plane with large normal component $\sin\theta_n$; the
refracted ray $\hat{\bm{k}}'_n$ (dashed grey) enters at the smaller
angle $\theta_t$ because of the high refractive index, and the
refracted TM direction $\hat{e}'_{p,n}$ (blue) is nearly tangential
($\sin\theta_t = O(1/|\ntilde|)$). Approximation~1
(\cref{prop:approx1}) replaces $\hat{e}'_{p,n}$ by $\hat{e}_{p,n}$;
the cross-term error is $O(1/|\ntilde|^2)$ because the small normal
component dominates the difference.}
\label{fig:incidence}
\end{figure}

Consider a planar interface between vacuum and a half-space of
biological tissue with complex refractive index $\ntilde =
\sqrt{\varepsilon_r - j\sigma/(\omega\varepsilon_0)}$. For skin at
28~GHz, $\ntilde \approx 4.49 - 1.79i$, with magnitude
$|\ntilde| \approx 4.83$. A plane wave incident at angle $\theta_n$
from the surface normal $\nhat$ admits the standard Fresnel TE and
TM amplitude coefficients $t_{s,n}$ and $t_{p,n}$, derived in
the supplementary information.

The incident wave $\psin = \psi_{s,n}\hat{e}_{s,n} +
\psi_{p,n}\hat{e}_{p,n}$ excites the transmitted field
$\EE^{\mathrm{trans}}_n$ inside the tissue, and the absorbed power
density at the surface, integrated over the skin depth, is
\begin{equation}\label{eq:Sab-single}
\Sab^{(n)}(\rr) = \frac{|x_{j(n)}|^2}{Z_0}\bigl(T_{s,n}|\psi_{s,n}|^2
+ T_{p,n}|\psi_{p,n}|^2\bigr)[\mu_n]_+,
\end{equation}
with $T_{s,n} = |t_{s,n}|^2 \cos\theta_n / \cos\theta_n$ the TE
power transmittance, $T_{p,n}$ the TM analogue, $Z_0 \approx 377\,\Omega$
the free-space impedance, $\mu_n = -\hat{\bm{k}}_n\cdot\nhat$ the
incidence cosine, and $[\mu_n]_+$ the back-face cull.

\subsection{The pseudo-Brewster collapse}

For materials with $|\ntilde| \geq 2 + \sqrt{3} \approx 3.73$, the
TE and TM power transmittances pinch together at all incidence
angles below the critical angle, a phenomenon documented in detail
in~\cite{WydaeghePB2024} and called the
\emph{pseudo-Brewster collapse}. Skin at mmWave satisfies the
condition: at 28~GHz both transmittances lie within $5\,\%$ of a
single scalar $\Tepskin \approx 0.54$ across
$\theta_n \in [0^{\circ}, 80^{\circ}]$. Substituting into
\cref{eq:Sab-single} collapses the polarisation structure to
$\Sab^{(n)} = \Sinc^{(n)}\,\Tepskin\,[\mu_n]_+$ with
$\Sinc^{(n)} = |x_{j(n)}|^2 |\psin|^2 / Z_0$, the geometric
absorption law on a lossy half-space~\cite{BornWolf}. The
polarisation-aware extension is invoked only for error bounds.

\subsection{The body as a discretised PO scatterer}

The user's body is the SMPL-X mesh of \cref{sec:arch}: a closed
surface $\Sigma(\thetab)$ tessellated into $T = 20\,908$ triangles.
For triangle $t$ at pose $\thetab$ we compute centroid
$\bm{c}_t(\thetab)$, outward unit normal $\nhat_t(\thetab)$, and
area $S_t(\thetab)$. We use a back-face cull
$V_n(t) = [-\hat{\bm{k}}_n \cdot \nhat_t]_+$ for incident direction
$\hat{\bm{k}}_n$, and a strict ray-mesh visibility test only
when fingertip-scale features matter (SI~\S{}S8).

The base-station path dictionary $\Jmat \in \Complex^{M \times N}$
encodes $N$ Sionna~RT-traced propagation paths from the BS to the
body's centroid region, each with per-element steering response
$\bm{a}_n \in \Complex^{M}$, arrival direction $\hat{\bm{k}}_n$,
polarisation $\psin$, and per-path amplitude $\beta_n$ (calibrated
against UL pilots; see SI~\S{}S5). For triangle $t$ and path $n$
the per-triangle incident wave at the body surface is
\begin{equation}\label{eq:incident-wave}
\psin \cdot \beta_n \cdot \bm{a}_n^{H}\xx \cdot
e^{-ik_0 \hat{\bm{k}}_n \cdot \bm{c}_t}.
\end{equation}

\subsection{The exposure operator and two enabling approximations}

The whole-body absorbed power $\Pabs(\thetab,\xx) = \int_{\Sigma}
\Sab\,dA$ becomes a Hermitian quadratic form
\begin{equation}\label{eq:Qabs}
\Pabs(\thetab,\xx) = \xx^{H} \Qabs(\thetab) \xx
\end{equation}
under two approximations whose error budgets are tight on biological
tissue at mmWave.

\begin{proposition}[TM polarisation direction; Approximation~1]\label{prop:approx1}
The physically exact transmitted TM field points along the refracted
direction $\hat{e}'_{p,n}$ inside the tissue. We replace
$\hat{e}'_{p,n}$ by the incident direction $\hat{e}_{p,n}$
(\cref{fig:incidence}). Self-terms are unaffected ($|\hat{e}_{p,n}|^2 =
|\hat{e}'_{p,n}|^2 = 1$); the error on TM--TM cross-terms is
$O(1/|\ntilde|^2) \lesssim 4\,\%$ for skin at 28~GHz. The error
shrinks at higher frequencies and on tissues with larger $|\ntilde|$.
\end{proposition}

\begin{proposition}[Universal depth coupling; Approximation~2]\label{prop:approx2}
The depth-integrated cross-coupling between paths $n$ and $n'$ is
$\Lambda_{nn'} = 1 / [\alpha_n + \alpha_{n'} - i(\beta_{n'} - \beta_n)]$
with $\alpha_n, \beta_n$ the per-path attenuation and phase rates
inside the tissue. Define the normalised coupling $\Gamma_{nn'} =
\Lambda_{nn'} / \sqrt{\Lambda_{nn}\Lambda_{n'n'}}$. We approximate
$\Gamma_{nn'} \approx 1$. For skin at 28~GHz over all incidence-angle
pairs $(\theta_n,\theta_{n'}) \in [0^{\circ},85^{\circ}]^2$ the error
satisfies $|1 - \Gamma_{nn'}| \leq 0.44\,\%$ with mean $0.17\,\%$;
the high refractive index compresses all refraction angles to within
$\sim\!10^{\circ}$ of normal, making $\alpha_n,\beta_n$ nearly
angle-independent.
\end{proposition}

The proofs are in the supplementary information. The combined effect of
\cref{prop:approx1,prop:approx2} is to factor the local absorbed
power as a squared norm:
\begin{equation}\label{eq:Sab-norm}
\Sab(\rr) = \bigl\|\tilde{\bm{G}}(\rr) \xx\bigr\|^2,
\end{equation}
with $\tilde{\bm{G}} \in \Complex^{3 \times M}$ the
\emph{exposure channel matrix}, defined per-element as
$\tilde{\bm{g}}_j(\rr) = \sum_{n: j(n)=j} \sqrt{\sigma/4\alpha_n}\,
\Fn(\rr)\psin\,e^{-ik_0 \hat{\bm{k}}_n \cdot \rr}$, with $\Fn$ the
per-path Fresnel transmission operator (rank-2, polarisation-aware,
defined formally in the supplementary information). The integral over the body
surface yields the operator
\begin{equation}\label{eq:Q-form}
\Qabs(\thetab) = \int_{\Sigma(\thetab)} \tilde{\bm{G}}(\rr)^{H}
\tilde{\bm{G}}(\rr) \,dA = \Jmat^{H}\Mmat(\thetab)\Jmat,
\end{equation}
where $\Mmat(\thetab) \in \Complex^{N\times N}$ is the body-and-scene
Gram and $\Jmat$ the array-hardware path dictionary. The
factorisation~\eqref{eq:Q-form} separates the per-slot precoder cost
from the per-pose body refresh, which is the structural property
that makes BS-side exposure-constrained beamforming
cheap~\cite{Hochwald2014,Ying2015}; we adopt the same factorisation
here for the SAR-reading branch.


% =====================================================================
\section{UE-anchored Kirchhoff render of the body channel}\label{sec:kirchhoff}

The per-triangle render geometry is shown in
\cref{fig:kirchhoff-schematic}. A representative path $n$ in the
BS path dictionary illuminates triangle $t$ on the body surface;
the local Fresnel reflection re-radiates a field whose contribution
at the phone $\bm{r}_p$ is the per-point integrand
$\mathcal{K}_n(\bm{r};\bm{r}_p)$. The body-mediated channel
coefficient is the integral of this kernel over the visible
surface, summed over BS paths.

\begin{figure}[!t]
\centering
\includegraphics[width=\columnwidth]{figures/fig_kirchhoff_schematic.pdf}
\caption{Per-triangle UE-anchored Kirchhoff render. The base station
illuminates triangle $t$ via path $n$ with arrival direction
$\hat{\bm{k}}_n$; the local Fresnel reflection re-radiates toward
the phone at $\bm{r}_p$ along the per-point arrival direction
$\hat{\bm{\eta}}(\bm{r};\bm{r}_p)$. The body-mediated channel
coefficient $\hbody(\bm{\theta},\bm{r}_p)$ is the integral of the
per-point kernel $\mathcal{K}_n$ over the visible surface, summed
over BS paths.}
\label{fig:kirchhoff-schematic}
\end{figure}

% =====================================================================

The body-mediated channel coefficient $\hbody(\thetab,\rr_p)$ is
obtained by computing, per surface point on $\Sigma(\thetab)$, the
field reflected by the local Fresnel reflection operator from the
incident BS multipath, and then propagating that re-radiated field
to the phone position $\rr_p$ via the free-space Green's function.
The result is a Kirchhoff radiation integral.

\subsection{The PO surface current}

For path $n$ incident on triangle $t$, the local reflected surface
current density is
\begin{equation}\label{eq:Js}
\bm{J}^{(n)}_t = 2\,\bigl(\nhat_t \times \HH^{(n)}_t\bigr)\,V_{nt},
\end{equation}
with $\HH^{(n)}_t$ the local incident magnetic field (derivable
from $\psin$ via the impedance $Z_0$) and $V_{nt} \in \{0,1\}$ the
visibility of triangle $t$ from path $n$'s arrival direction
(back-face cull plus optional ray-mesh self-shadow). For the
phone-on-body geometries that motivate the work, the back-face cull
suffices; the strict visibility test is invoked only for fingertip
features and is detailed in the supplementary information.

\subsection{Propagation to the phone}

The phone is at world position $\rr_p$. The free-space Green's
function from a surface point $\rr$ to the phone is
\begin{equation}
G(\rr,\rr_p) = \frac{e^{jk_0 \|\rr-\rr_p\|}}{4\pi|\rr-\rr_p|},
\end{equation}
and the per-point arrival direction at the phone is $\etahat(\rr;\rr_p)
= (\rr_p - \rr) / |\rr_p - \rr|$. By the Kirchhoff integral, the
field arriving at the phone from the body's reflected current is
\begin{equation}\label{eq:kirchhoff-integral}
\hbody(\thetab,\rr_p) = \sum_{j=1}^{M} \bigl[\bm{a}_n\bigr]_j
\sum_{n=1}^{N} \int_{\Sigma_n^{+}(\thetab)} \mathcal{K}_n(\rr;\rr_p)\,dA,
\end{equation}
with the per-point integrand
\begin{equation}\label{eq:kernel}
\mathcal{K}_n(\rr;\rr_p) = K\,\bm{g}_{\mathrm{UE}}(\etahat)^{H}\psin
\,G(\rr,\rr_p)\,e^{-jk_0\hat{\bm{k}}_n\cdot\rr} \,,
\end{equation}
where $K = jk_0 r_0$ is the pseudo-Brewster reflection amplitude
($r_0 = 1 - \Tepskin \approx 0.46$) and
$\bm{g}_{\mathrm{UE}}(\etahat)$ is the phone's antenna pattern in
the arrival direction (isotropic by assumption: $|\bm{g}_{\mathrm{UE}}|^2
\equiv 1$). The visible-region $\Sigma_n^{+}$ is the subset of
$\Sigma$ that satisfies the back-face cull for path $n$ and is in
line of sight of $\rr_p$.

\subsection{Discretisation and the Kirchhoff operator}

Discretising the integral over the $T_{\mathrm{vis}}$ visible
triangles yields the matrix form
\begin{equation}\label{eq:Lambda-op}
\hbody = \mathbf{a}_{\mathrm{LOS}}^{H}\Lambda_{\mathrm{KH}}(\thetab,\rr_p)\,
\bm{1}_{T_{\mathrm{vis}}},
\end{equation}
where $\Lambda_{\mathrm{KH}} \in \Complex^{M \times T_{\mathrm{vis}}}$
is the \emph{body-to-phone Kirchhoff operator} with entries
$\Lambda_{\mathrm{KH}}[m,t] = \mathcal{K}_n(\bm{c}_t;\rr_p) \cdot
\beta_n \cdot S_t \cdot \bigl[\bm{a}_n\bigr]_m$ summed over the
paths $n$ landing on triangle $t$. The operator $\Lambda_{\mathrm{KH}}$
carries the full pose-dependence of the body-mediated channel; its
singular spectrum, analysed in \cref{sec:closed-loop}, governs the
calibration identifiability.

\begin{remark}[Translation phasor identity]
For a small displacement $\bm{t}$ of the body's centroid (e.g., a
walking step), the operator transforms as $\Lambda_{\mathrm{KH}}(\rr_0
+ \bm{t}) = \Phimat(\bm{t})^{H}\Lambda_{\mathrm{KH}}(\rr_0)$ with
$\Phimat$ a per-element diagonal phasor sandwich. This is the
analogue of the standard
array-translation phase identity and lets the per-slot operator
refresh under translation cost only one diagonal multiply.
\end{remark}

The phone-on-body geometry has exactly one receiver --- the phone,
about $50\,$cm from the body --- so the backward Kirchhoff render
that queries each visible triangle and integrates the local
Fresnel-reflected field anchored at the phone is the appropriate
primitive. We measure $\sim\!100\,$ms per pose for the Kirchhoff
render on a single CPU core at $T = 20\,908$, $T_{\mathrm{vis}}
\sim 10\,000$, $M = 64$, on the SMPL-X body of \cref{sec:arch}.


% =====================================================================
\section{Closed-loop calibration: UL pilots and \texorpdfstring{$\gammab$}{gamma}}\label{sec:closed-loop}
% =====================================================================

The Kirchhoff render of \cref{sec:kirchhoff} is a \emph{prediction}
of $\hbody$ from the on-device twin. The actual body-mediated channel
on the air differs from the prediction because the SMPL-X mesh is
not the user's exact anatomy, the dielectric assumption $\Tepskin
\approx 0.54$ is a population mean, and the IMU pose is not the
ground-truth pose. We call this aggregate the \emph{twin gap}.
Closing it requires a regression of measured against predicted
channel, fit per-slot or per-second over the cooperative downlink
in which the user is enrolled.

\subsection{Residual decomposition}

Let $\hh^{\mathrm{meas}}$ be the channel observed at the BS via
uplink pilots and $\hh^{\mathrm{pred}}(\thetab) =
\hBS + \hbody^{\mathrm{pred}}(\thetab)$ be the cascaded channel
predicted by the twin. The residual $\Delta = \hh^{\mathrm{meas}} -
\hh^{\mathrm{pred}}$ decomposes into BS-side and body-side
components:
\begin{equation}\label{eq:residual}
\Delta = \underbrace{\Phimat(\rr_p)\Jmat\,\delta\betab}_{\text{BS-side}}
+ \underbrace{\Lambda_{\mathrm{KH}}(\thetab,\rr_p)\,\delta\bm{q}}_{\text{body-side}}
+ \bm{n}_{\mathrm{p}},
\end{equation}
where $\delta\betab$ are the per-path amplitude corrections
(BS-side, calibrated against uplink pilots; see SI~\S{}S5),
$\delta\bm{q}$ is a
per-triangle correction factor on the body-mediated current, and
$\bm{n}_{\mathrm{p}}$ is the pilot noise. The two sides are
identifiable from the same uplink pilot stream as long as the
column spaces of $\Phimat\Jmat$ and $\Lambda_{\mathrm{KH}}$ are
distinguishable.

\subsection{Body-side parametrisation: a $K$-mode SVD basis}

The body-side correction $\delta\bm{q}$ has $T_{\mathrm{vis}} \sim
10^{4}$ entries, far more than the $M = 64$ pilot equations admit.
We project $\delta\bm{q}$ onto a $K$-dimensional basis
$\bm{V}_K \in \Complex^{T_{\mathrm{vis}}\times K}$, write
$\delta\bm{q} \approx \bm{V}_K \gammab$, and fit the body-side
calibration coefficient $\gammab \in \Complex^{K}$ by ridge least
squares against the residual. The basis $\bm{V}_K$ is the leading
$K$ right singular vectors of $\Lambda_{\mathrm{KH}}$ itself, with
$K = K_{99}$ chosen adaptively as the smallest $K$ that captures
$99\,\%$ of the operator's spectral mass. The operator-spectrum
sweep of SI~\S{}S7 finds $K_{99} \in \{1, 2, 4\}$ across the BS
geometries we consider.

\begin{proposition}[Identifiability]\label{prop:ident}
Let $\bm{U}_K$ be the leading-$K$ left singular vectors of
$\Lambda_{\mathrm{KH}}$, with corresponding singular values $\sigma_k
> 0$. The body-side ridge least squares with basis $\bm{V}_K$ equal
to the corresponding right singular vectors is exactly
identifiable when the residual lies in $\mathrm{span}(\bm{U}_K)$.
The fit is non-trivial as long as $K \leq M$ and the pilot SNR
exceeds $K \sigma_K^{-2}$.
\end{proposition}

The proof is a standard truncated SVD argument and is given in
the supplementary information.

\subsection{Calibration robustness to IMU drift}

The body-side $\gammab$ absorbs both pose error (residual IMU drift
after attitude filtering) and dielectric mismatch in the same
regression. The mechanism is the same: both error sources produce
residuals in the column space of $\Lambda_{\mathrm{KH}}$ at first
order, so the same $K$-mode SVD basis that captures the dominant
body-to-phone modes also captures the dominant residual directions.
The closed loop is therefore robust to consumer-grade IMU drift in
the pose telemetry: the calibration step recovers the predicted
channel from a noisy pose input by absorbing the residual into the
operator's leading modes. The IMU-noise sensitivity sweep (SI~\S{}S6)
quantifies the rate cost of this drift across
$\sigma_{\mathrm{joint}} \in \{0, 2, 4, 8, 16\}^{\circ}$.


% =====================================================================
\section{Pose-differentiable rate and human-in-the-loop control}\label{sec:control}
% =====================================================================

\subsection{The cascaded SINR and the rate}

For an $M$-element BS using maximum-ratio transmission $\xx =
\sqrt{P_{\mathrm{tx}}}\,\hh / \|\hh\|$, the post-precoding SINR at
the phone is $\mathrm{SINR}(\thetab) = P_{\mathrm{tx}} \|\hh(\thetab)\|^{2}
/ N_0$. The achievable rate per Hz is the Shannon expression
\begin{equation}\label{eq:rate}
R(\thetab) = \log_2\bigl(1 + \mathrm{SINR}(\thetab)\bigr)\,.
\end{equation}
The cascaded structure of \cref{eq:cascaded} carries the full
pose-dependence: $\nabla_{\thetab} R$ is non-zero everywhere
$\mathrm{SINR}>0$, with magnitude controlled by
$\partial R / \partial \mathrm{SINR} = (\ln 2)^{-1}/(1 + \mathrm{SINR})$.

\subsection{Pose-space restriction and the latent gradient}\label{sec:vposer-latent}

The axis-angle vector $\thetab \in \Reals^{66}$ is over-parametrised
for our purposes. Most of $\Reals^{66}$ corresponds to anatomically
impossible configurations (joints outside their range of motion,
incompatible torso--limb couplings). Optimisation on $\Reals^{66}$
would routinely propose contortionist poses that the user cannot
actuate. We therefore restrict the optimisation to the manifold of
plausible body poses, which we parametrise by the VPoser variational
prior~\cite{Pavlakos2019VPoser}. VPoser is a VAE trained on the AMASS
motion-capture corpus~\cite{AMASS2019}; its decoder
$\bm{G}: \Reals^{32} \to \Reals^{63}$ maps a 32-dimensional latent
$\bm{z}$ to the 21 SMPL-X body-pose joints (excluding the global
orientation and the hands), and the prior $p(\bm{z}) = \mathcal{N}(\bm{0},
\bm{I})$ assigns high probability to natural poses. Any latent
$\bm{z}$ decodes to an anatomically valid configuration by
construction.

The pose-differentiable rate becomes a function of the latent:
$\thetab(\bm{z}) = \bm{G}(\bm{z})$,
$\hh(\bm{z}) = \hBS + \hbody(\thetab(\bm{z}))$, and
$R(\bm{z}) = \log_2(1 + P_{\mathrm{tx}}\|\hh(\bm{z})\|^{2}/N_0)$.
The full chain
\begin{equation}\label{eq:autodiff-chain}
\bm{z} \xrightarrow{\bm{G}} \thetab \xrightarrow{\mathrm{SMPL\text{-}X}}
\Sigma \xrightarrow{\mathrm{Kirchhoff}} \hbody \xrightarrow{}
\mathrm{SINR} \xrightarrow{\log_2(1 + \cdot)} R
\end{equation}
is a composition of differentiable maps. The Kirchhoff render
$\Sigma \mapsto \hbody$ is linear in the per-triangle vertex
positions so its Jacobian is trivial; SMPL-X is a sequence of
linear-blend-skinning matrix products with closed-form derivatives;
the VPoser decoder is a feed-forward MLP with a continuous
6-D rotation representation. We compute $\nabla_{\bm{z}} R$ end-to-end
via reverse-mode autodifferentiation, requiring one forward pass and
one backward pass per iteration (no finite-difference oracle on the
underlying physics).

Comfort is a naturalness constraint in latent space:
$C(\bm{z}) = \|\bm{z}\|^{2} \leq \rho_\mathrm{nat}^{2}$, anchored
at the VPoser prior origin (not at the user's current
$\bm{z}_0$).~Under the VAE prior $p(\bm{z}) =
\mathcal{N}(\bm{0},\bm{I})$~\cite{Pavlakos2019VPoser}, the
constraint $\|\bm{z}\| \leq \rho_\mathrm{nat}$ means the prior
log-probability of every admissible pose is above
$\log p(\bm{0}) - \rho_\mathrm{nat}^{2}/2$. We use
$\rho_\mathrm{nat} = 2.0$, which keeps poses within $2$ nats of the
prior mode --- broadly natural, but excluding the
contortionist tail (yoga splits, gymnast bridges) that VPoser
assigns small density. The user's current pose $\bm{z}_0$ enters
only as the SINR baseline against which gain is measured.

\subsection{The human-in-the-loop control law}

The latent gradient $\nabla_{\bm{z}} R$ is computed on the device,
not actuated mechanically. The user is the actuator. Rather than
issue a one-shot pose target $\thetab^{*}$, which would require the
BS to know the user's exact joint capabilities and current intent,
the BS issues at each tick a small latent step $\Delta\bm{z}$ along
$\nabla_{\bm{z}} R$ projected onto the naturalness ball $\|\bm{z}\| \leq \rho_\mathrm{nat}$, and the SMPL-X
decoder turns $\Delta\bm{z}$ into the corresponding skeletal
configuration shown to the user as a haptic or visual cue. The
user accepts or rejects; the next tick observes the new IMU pose,
re-encodes it to a new $\bm{z}_0$, runs a fresh $\gammab$ fit on
the new uplink pilot, and emits the next step
(\cref{alg:closed-loop}). The loop is projected gradient ascent on
a slowly-drifting non-stationary objective with measurement
validation at each iterate.

\begin{algorithm}[!t]
\caption{$K$-iteration latent projected gradient ascent.}
\label{alg:closed-loop}
\begin{algorithmic}[1]
\State \textbf{Input:} baseline pose $\thetab^{(0)}$, naturalness
radius $\rho_\mathrm{nat}$, step size $\eta$, iteration budget $K$
\State Encode $\bm{z}_0 \leftarrow \bm{G}^{-1}(\thetab^{(0)})$
via the VPoser encoder
\State Initialise $\bm{z}^{(0)} = \bm{z}_0$ if $\|\bm{z}_0\| \leq \rho_\mathrm{nat}$ else
$\bm{z}^{(0)} = \bm{z}_0\,\rho_\mathrm{nat}/\|\bm{z}_0\|$
\For{$k = 0, 1, \ldots, K-1$}
\State Decode $\thetab^{(k)} \leftarrow \bm{G}(\bm{z}^{(k)})$ and run
SMPL-X forward map $\thetab^{(k)} \mapsto \Sigma(\thetab^{(k)})$
\State Render $\hbody^{(k)} \leftarrow$ UE-anchored Kirchhoff
integral on $\Sigma(\thetab^{(k)})$
\State Observe UL pilot residual $\Delta^{(k)} = \hh^{\mathrm{meas}}
- (\hBS + \hbody^{(k)})$
\State Fit $\gammab^{(k)}$ over the $K_{99}$-mode SVD basis of
$\Lambda_{\mathrm{KH}}^{(k)}$
\State Update predicted channel
$\hat\hh^{(k)} = \hBS + \hbody^{(k)} +
\Lambda_{\mathrm{KH}}^{(k)} \bm{V}_{K_{99}} \gammab^{(k)}$
\State Compute latent gradient $\bm{g}^{(k)} = \nabla_{\bm{z}}
R(\bm{z}^{(k)})$ by reverse-mode autodifferentiation through
\cref{eq:autodiff-chain}
\State Project step onto the naturalness ball:
$\Delta\bm{z}^{(k)} = \mathrm{Proj}_{\|\bm{z}\|\leq \rho_\mathrm{nat}}\!
\bigl(\bm{z}^{(k)} + \eta\,\bm{g}^{(k)}\bigr) - \bm{z}^{(k)}$
\State User actuates: $\bm{z}^{(k+1)} = \bm{z}^{(k)} + \Delta\bm{z}^{(k)}$
\EndFor
\State \Return $\thetab^{(K)} = \bm{G}(\bm{z}^{(K)})$,
$R^{(K)} = R(\bm{z}^{(K)})$
\end{algorithmic}
\end{algorithm}

The user need not see the gradient itself. A simple haptic or
cue (a torso-yaw bar or a small arrow on the screen) suffices
because the actuation is body-scale and the rate signal refreshes at
$\sim 10\,$Hz. This is one of the two semantic outputs the
network's wireless control loop issues to the user; the other is
the live exposure reading of \cref{sec:saridx}.


% =====================================================================
\section{Numerical study: existence, reachability, failure modes}\label{sec:numerical}
% =====================================================================

\subsection{Existence, reachability, failure modes}

The closed loop of \cref{sec:closed-loop,sec:control} either reaches
a useful operating point or it does not. Three nested questions
characterise when:
\begin{enumerate}
\item[\textbf{(E)}] \textbf{Existence.} For a given (BS-geometry,
scene, baseline-pose) tuple, does \emph{any} comfortable latent
move $\bm{z}$ with $\|\bm{z}\| \leq \rho_\mathrm{nat}$ achieve a
target rate gain $R(\bm{z}) - R(\bm{z}_0) \geq \Delta R_{\mathrm{tgt}}$?
\item[\textbf{(R)}] \textbf{Reachability.} Conditional on existence,
does $K$-iteration UL-pilot-corrected projected gradient ascent
on $\bm{z}$ from the baseline reach within 1~dB of the comfortable
optimum?
\item[\textbf{(F)}] \textbf{Failure modes.} For trajectories that
fail to reach, what fails: stagnant ascent, local minimum trap,
calibration drift, or IMU-noise floor?
\end{enumerate}

\subsection{Setup}

All experiments use the SMPL-X neutral mesh with baseline poses
sampled from $50$ AMASS walks at $30\,$fps, yielding $\sim\!13\,000$
plausible pose frames; each baseline frame $\thetab^{0}$ is encoded
to a VPoser latent $\bm{z}_0$ that we use as the SINR baseline (not
as a constraint anchor). The naturalness ball
$\|\bm{z}\| \leq \rho_\mathrm{nat} = 2.0$ keeps every admissible
pose within $2$ nats of the VPoser prior mode, ruling out the
contortionist tail.

The base station is an $8\times 8$ URA at $28\,$GHz with
$\lambda/2$ spacing on a $35\,$m mast, transmitting $43\,$dBm with
a noise floor of $-94\,$dBm. The panel is downtilted and yawed to
broadside the served plaza (panel-normal vector
$\hat{\bm{n}}_\mathrm{BS} = (0.30, 0.86, -0.41)$, pointing at the
ground centroid of the served LOS region at $(33, 91, 1.5)\,$m). The user is at
$1.2\,$m height; the phone is held at chest height, $40\,$cm in
front of the body centroid. Rates are reported as plain Shannon
$\log_2(1 + \mathrm{SINR})$ (no MCS cap).

\paragraph*{On-device pose estimation.} The on-device twin
consumes a per-tick SMPL-X pose estimate $\hat{\bm{\theta}}$
produced by a sparse-IMU body-pose inference pipeline running
on the user's consumer devices. Recent works in this lineage --
DIP~\cite{DIP2018}, TransPose~\cite{TransPose2021},
IMUPoser~\cite{Mollyn2023IMUPoser}, MobilePoser~\cite{Xu2024MobilePoser}
-- regress full-body pose from sparse inertial readings at
real-time rates on a phone. State-of-the-art consumer-only
accuracy is $\sim\!24^{\circ}$ mean per-joint rotation error
(MPJRE) at $60\,$fps on a current smartphone~\cite{Xu2024MobilePoser};
research-grade body-suit estimators (e.g., PIP at $6$ IMUs)
reach $\sim\!12^{\circ}$. The IMU sensitivity sweep
(SI~\S{}S6) characterises the closed-loop SINR cost of
operating on $\hat{\bm{\theta}}$ instead of the noiseless
$\bm{\theta}$ across a per-joint attitude noise range
$\sigma_\mathrm{joint} \in \{0, 2, 4, 8, 16\}^{\circ}$ that
spans research-grade body suits ($\sigma \approx 4^{\circ}$)
through current consumer-only operating points
($\sigma \approx 16^{\circ}$). The peak-APD location
sensitivity sweep (SI~\S{}S7) does the same for the cohort-side
absorbed-power reading.

The propagation environment is the canonical Munich scene
distributed with Sionna~RT~\cite{Hoydis2023Sionna}, a high-fidelity
city-block model widely used as a benchmark for ray-tracing studies.
We pick $18$ Rx positions ($9$ LOS, $9$ NLOS;
\cref{fig:munich-scene}) interactively from a 3D rendering of the
Sionna geometry: LOS positions sit on the open plaza visible from
the panel broadside, NLOS positions sit on side streets shadowed
by intervening buildings. Each Rx is then traced with Sionna RT
(maximum reflection depth $5$, specular reflection $+$ refraction
$+$ edge diffraction enabled, $10^{7}$ samples per source); the
chosen LOS cells return $\sim\!500$ paths, the NLOS cells $\sim\!300$. We do not need millimetre-accurate scene
geometry: the BS-side $\bm{\beta}$ calibration of
\cref{sec:closed-loop} absorbs the per-path amplitude and phase
discrepancies between the simulated scene and the live channel,
provided the ray-traced angles-of-arrival are approximately right.
Per-element steering across the URA is applied analytically via
planewave phase ramps in the per-path direction of departure from
the array centroid. The body-mediated channel $\hbody$ is then
computed by feeding the per-path body-incident bundle to the
Kirchhoff render of \cref{sec:kirchhoff}.

\begin{figure}[!t]
\centering
\includegraphics[width=\columnwidth]{rx_grid_annotated.png}
\caption{Munich scene (Sionna RT canonical) with the BS at
$(8.5, 21.7, 35)\,$m mast height; yellow arrow shows the panel
broadside direction. The $9$ LOS Rx (green) lie on the open
plaza and side avenues visible from the panel; the $9$ NLOS Rx
(red) sit on streets where intervening buildings shadow the BS.
Rx selected by hand from a 3D rendering of the Sionna geometry.
Phone height $1.5\,$m, body centroid $0.40\,$m behind phone
facing the BS.}
\label{fig:munich-scene}
\end{figure}

\subsection{Hero example}

\Cref{fig:hero-pose} illustrates the central effect on the
highest-gain Munich NLOS scene (scene~$17$, $\mathrm{Rx} =
(-69, 115, 1.5)\,$m on a deeply-shadowed alley north-west of
the BS, scene loss $-109\,$dB). A latent perturbation within the
naturalness ball $\|\bm{z}\|\leq\rho_\mathrm{nat}=2.0$ produces
an $8.6\,$dB cascaded-SINR difference between the
encoder-projected baseline pose ($10.9\,$dB) and the best of
$N=100$ random latent samples ($19.5\,$dB). With the BS-side
bundle starved by the deep shadow, the body becomes the dominant
scattering term and the pose move acts as a body-RIS rescue.
Panels~(a) and~(b) show the body silhouettes from the BS
perspective with per-triangle BS-illumination shading; panel~(c)
plots the SINR of each random latent sample, with baseline and
best marked.

\begin{figure*}[!t]
\centering
\includegraphics[width=0.95\textwidth]{fig_hero_munich.pdf}
\caption{Hero example, Munich NLOS scene~$17$ (Rx at
$(-69, 115)\,$m on a deeply-shadowed alley, scene loss
$-109\,$dB). (a)~Baseline pose (the VPoser encoder mean of an
AMASS walk frame), cascaded SINR $10.9\,$dB. (b)~Best of $N=100$
random latent samples within the naturalness ball, $19.5\,$dB.
The shading colours each visible triangle by its BS-illumination
cosine. (c)~Cascaded SINR of all $100$ random latent samples;
the baseline (red dashed) and best-found (green dashed) are
marked, with $\Delta = 8.6\,$dB.}
\label{fig:hero-pose}
\end{figure*}

\subsection{Existence}

\Cref{fig:existence} reports the per-scene SINR gain achieved by a
naturalness-bounded latent perturbation over the $18$ Munich
scenes (\cref{fig:munich-scene}) stratified by LOS/NLOS. The nine
LOS scenes carry $53$~to $58\,$dB of SINR from the BS multipath
bundle and admit small uplifts: gains span $0.6$~to $2.4\,$dB
(mean $1.7\,$dB). The nine NLOS scenes split into three regimes
by scene loss. Mild-shadow NLOS ($\sim\!75$~to $90\,$dB scene
loss) returns the same $1$~to $4\,$dB pose uplift as LOS. Strong-shadow
NLOS ($\sim\!90$~to $100\,$dB) opens $4$~to $8\,$dB of pose
leverage; the body now contributes a non-trivial scattering term.
The deep-shadow case (scene~17, $-69\,$m, $115\,$m, scene loss
$-109\,$dB) is the body-RIS regime: the BS-side bundle alone
delivers only $11\,$dB of SINR; a naturalness-bounded pose move
adds $\bm{8.6\,}$dB, with the body acting as the dominant
scatterer back to the phone.

\begin{figure}[!t]
\centering
\includegraphics[width=\columnwidth]{fig_munich_perscene.pdf}
\caption{Per-scene cascaded SINR over the $18$ Munich scenes
($9$ LOS, $9$ NLOS). Light bar: baseline pose $\bm{z}_0$
(VPoser-encoded AMASS walk frame). Dark bar: best of $N=100$
random latent samples within the naturalness ball
$\|\bm{z}\| \leq \rho_\mathrm{nat} = 2.0$~\cite{Pavlakos2019VPoser}.
The deep-shadow scene~17 (rightmost NLOS) shows the body-RIS
rescue regime: $+8.6\,$dB of pose leverage where the BS-side
bundle alone is starved.}
\label{fig:existence}
\end{figure}

\subsection{Reachability}

\Cref{fig:reachability} reports the reachability of the
projected latent gradient-ascent loop of \cref{alg:closed-loop} at
$K \in \{2, 5, 10, 15, 20\}$ iterations. ``Success'' is defined as
terminal SINR within $1\,$dB of the brute-force optimum (the
naturalness-ball maximum found by $N=100$ random samples). At
$K = 20$, the loop reaches the optimum on $8/9$ LOS scenes and
$8/9$ NLOS scenes. The two failures are dominated by local
plateaus in the latent-space rate landscape, including the
deep-shadow scene~17 where the loop terminates within $1.1\,$dB
of the brute-force ball optimum. The IMU-noise sensitivity sweep
(SI~\S{}S6) shows that the rate loss the loop pays under
consumer-grade IMU operating points is small.

\begin{figure}[!t]
\centering
\includegraphics[width=\columnwidth]{fig_munich_reachability.pdf}
\caption{Reachability fraction over $18$ Munich scenes. The
closed loop of \cref{alg:closed-loop} is run for
$K \in \{2, 5, 10, 15, 20\}$ iterations of autodiff projected
gradient ascent on the VPoser latent. Success is terminal SINR
within $1\,$dB of the brute-force naturalness-ball optimum.}
\label{fig:reachability}
\end{figure}

\subsection{Per-joint sensitivity}

\Cref{fig:per-joint} ranks the joint-space sensitivity
$|\partial R / \partial \theta_j|$ across the 22 SMPL-X body
joints, averaged over $50$ AMASS walk poses at the same canonical
scene. Although the closed loop optimises in latent space (not
joint space), the joint-space ranking is interpretable feedback
for the user: the top five joints are the pelvis, the two hips,
and the lower two spine joints, and they together carry over
$80\,\%$ of the gradient mass. Wrists, elbows, ankles, and the
head are essentially insensitive: they do not move the body-to-phone
Kirchhoff coupling under normal phone-holding geometries. The
VPoser decoder respects this anatomy automatically, so the latent
gradient $\nabla_{\bm{z}} R$ correspondingly concentrates on the
modes that actuate torso joints rather than periphery joints.

\begin{figure}[!t]
\centering
\includegraphics[width=\columnwidth]{fig_per_joint_sensitivity.pdf}
\caption{Per-joint pose-rate sensitivity $|\partial R / \partial
\theta_j|$ ranked across the 22 SMPL-X body joints, averaged
$\pm$IQR over $50$ AMASS walk poses. The top five joints (pelvis,
left/right hip, spine1, spine2) carry the bulk of the gradient
mass; the anatomical inset highlights their location on the body.
Periphery joints (wrists, elbows, ankles, head) are essentially
insensitive.}
\label{fig:per-joint}
\end{figure}

\subsection{Comfort-vs-rate Pareto}

\Cref{fig:pareto} reports the SINR gain achievable as a function
of the latent comfort radius $\|\bm{z} - \bm{z}_0\|$, sampled with
$500$ random latent perturbations across the ball
$\|\bm{z} - \bm{z}_0\| \leq \rho_z = 1.5$ at one representative
NLOS-class scene from the parametric scalar-loss companion sweep
(see SI~\S{}S3). The Pareto envelope is monotone in
$\|\bm{z}-\bm{z}_0\|$: a $0.5$-radius latent step buys
$\sim 6\,$dB of SINR gain, a unit step buys $\sim 7.5\,$dB, and the
marginal return saturates above $\|\bm{z}-\bm{z}_0\| \approx 1.2$.
The cloud of below-frontier samples shows that random
comfort-bounded latent moves can equally lose SINR; the directed
projected gradient ascent of \cref{alg:closed-loop} is what tracks
the frontier.

\begin{figure}[!t]
\centering
\includegraphics[width=\columnwidth]{fig_comfort_pareto.pdf}
\caption{Comfort-vs-rate Pareto envelope at one representative
NLOS-class scene (parametric scalar-loss companion sweep, SI~\S{}S3).
Random latent perturbations ($N=500$) fill the cloud; the Pareto
frontier (red) is the upper envelope. A unit-radius latent step
buys $\sim 7.5\,$dB of SINR gain.}
\label{fig:pareto}
\end{figure}

\subsection{Comparison against baselines}

\Cref{fig:baselines} compares four precoder families across the
three scalar-loss bands of the parametric companion sweep
(SI~\S{}S3): no-twin zero-forcing, no-twin maximum-ratio
transmission, T-pose-prior twin (BS assumes the canonical SMPL-X
T-pose), and the full closed-loop RIHB. In the LOS-attenuated
band all twin-aware methods reach the same operating point;
only no-twin ZF degrades because it inverts against the wrong
channel ($-18\,\%$ relative to no-twin MRT). In the body-comparable
band ($55\,$dB scalar loss), the RIHB closed loop gains $7.1\,\%$
over no-twin MRT ($48\,$Mbps) while the T-pose-prior twin gains
$5.8\,\%$. In the body-dominant band ($80\,$dB scalar loss), both
T-pose and RIHB rescue the link ($+13.8\,\%$ over no-twin MRT). The
twin's value depends on the operating regime: the body-comparable
band is where the body-side calibration most clearly separates RIHB
from a static-pose prior.

\begin{figure}[!t]
\centering
\includegraphics[width=\columnwidth]{fig_baselines.pdf}
\caption{Mean rate by scalar-loss band and precoder family on $30$
AMASS scenes (10 per band) from the parametric companion sweep
(SI~\S{}S3). Error bars are $\pm$one standard deviation
across scenes. No-twin ZF degrades in LOS-attenuated. RIHB
strictly wins in the body-comparable band where the body-side
calibration is decisive.}
\label{fig:baselines}
\end{figure}


% =====================================================================
\section{Live SAR reading and connection to RF-EMF cohort studies}\label{sec:saridx}
% =====================================================================

\subsection{The absorbed-power reading}

At each twin tick, the on-device twin evaluates \cref{eq:Sab-norm}
on the visible triangles for the current precoder $\xx$ and
integrates to obtain the whole-body absorbed power $\Pabs$ and
the cumulative dose $\int \Pabs\,dt$ over the user's session.
The reading is displayed in milliwatts (instantaneous) and
milli-watt-hours (cumulative). At the canonical operating point
($43\,$dBm BS power, $30\,$m range, $28\,$GHz, $8\times 8$ URA),
$\Pabs$ sits at $0.58\,\%$ of the ICNIRP $4\,$W/m$^{2}$
general-public limit (the supplementary information).
Compliance is therefore not the binding constraint on $\xx$ at
the FR2 operating points considered here; the binding constraint
is the link budget itself. The reading nevertheless serves a
non-engineering function: it puts the user, the regulator, and
the epidemiologist on the same page about what the body is
actually absorbing at every moment.

The reading reported here is the whole-body integrated $\Pabs$
and its time integral, both of which are robust to the on-device
pose estimate (SI~\S{}S7): at the consumer-grade IMU operating
point ($\sigma_\mathrm{joint} = 16^{\circ}$), the on-device
$\Pabs$ tracks the true value to within a ${\sim}3\,\%$ mean bias
and ${\pm}4\,\%$ standard deviation across the $18$ scenes; at
$\sigma_\mathrm{joint} = 4^{\circ}$ (research-grade IMU suit),
the reading is within ${\pm}1\,\%$. The peak-local $\Sab$ over
the ICNIRP $4\,$cm$^{2}$ surface-average window is not reported
per-tick: the underlying argmax over body triangles is unstable
under realistic IMU pose noise (peak location standard deviation
$\sim\!50\,$cm; SI~\S{}S7), so per-tick anatomical resolution of
the hot-spot is out of scope. Time-aggregated or larger-window
peak-$\Sab$ statistics may recover a stable localised reading;
we leave that to follow-up work.

\subsection{One operator, two semantic outputs}

The reading is computed on the same per-triangle Kirchhoff machinery
as $\hbody$ but with the integrand evaluated in absorption mode
($\Sab = \|\tilde{\bm{G}}\xx\|^{2}$) rather than reflection mode
($\hbody = $ Kirchhoff radiation integral). The two modes share the
same operator factorisation \cref{eq:Q-form} and the same per-element
exposure channel matrix $\tilde{\bm{G}}$. The SAR-matrix
machinery~\cite{Hochwald2014,Ying2015,Ying2017} that supplies
$\Qabs$ to the BS-side exposure-constrained
beamformer~\cite{Hochwald2014,Ying2015} is the same machinery that supplies
$\Sab(\rr)$ to the user-facing reading. The reading is the
user-facing semantic output of an object the network already
maintains.

\subsection{Aggregation for cohort exposure studies}

Prospective RF-EMF cohort studies (the COSMOS and HERMES
prospective cohorts~\cite{Roosli2010,Roosli2021}, the European
CLUE-H consortium\footnote{CLUE-H: cohort linkage of mobile-phone
exposure and health outcomes; consortium of EU prospective cohort
studies aiming at high-quality bio-outcome assessment with
measurement-grade exposure side. Confirm canonical reference at
submission.}) measure bio/health outcomes in well-characterised
populations and link them to RF-EMF exposure. The exposure side
currently rests on self-reported phone-use questionnaires
plus modelled SAR proxies. Per-individual measured absorbed-energy
time series, opt-in-aggregated from the reading, would upgrade the
exposure side from self-report to direct measurement.

The reading's privacy delta is small: only the time series of total
absorbed power (a scalar per twin tick) need leave the device for
cohort use, with no body-twin parameters and no surface deposition
maps. Aggregation cadence is the user's choice; longitudinal cohort
studies need monthly or quarterly summaries, not per-slot streaming.

The CLUE-H consortium~\cite{CLUEH2025} is a coordinated set of EU
prospective cohort studies aimed at high-resolution biological /
health-outcome assessment with measurement-grade RF-EMF exposure
characterisation, the consortium-scale extension of single-cohort
studies in the lineage of COSMOS~\cite{Roosli2021},
HERMES~\cite{Roosli2010}, and the methodology work of van~Wel and
collaborators~\cite{vanWel2024}. The consortium's design currently
binds on the exposure side: questionnaire-reported phone use and
modelled SAR proxies are coarse compared to the cohort outcomes
they are linked against. Per-individual reading-aggregated
absorbed-energy time series upgrade the exposure side from
self-report to direct measurement and are obtained without any
additional in-cohort instrumentation cost. \Cref{fig:sab-reading}
illustrates the on-device computation: per-triangle
absorbed-power-density on the SMPL-X body, an instantaneous
$\Pabs$ time series across a $30\,$s walk, and the $100$-user
cumulative-dose distribution showing that the operating point
sits well below the regulatory ceiling.

\begin{figure*}[!t]
\centering
\includegraphics[width=0.95\textwidth]{fig_sab_receipt.pdf}
\caption{Live exposure reading on the SMPL-X body twin.
(a)~Per-triangle $\Sab$ heatmap on the body surface for one
representative pose; the cyan star marks the phone position.
(b)~Instantaneous $\Pabs$ over a $30\,$s AMASS-walk trace, with
cumulative dose on the secondary axis. (c)~ECDF of $30\,$s
cumulative dose across $100$ random users; the orange dashed line
marks the $4\,$W/m$^{2}$ ICNIRP-equivalent reference, well to the
right of the user distribution.}
\label{fig:sab-reading}
\end{figure*}


% =====================================================================
\section{Discussion}\label{sec:discussion}
% =====================================================================

\subsection{Bandwidth and per-slot compute}\label{sec:bandwidth}

\Cref{tab:cadence} lists the refresh cadences and per-update
wall-clock times for each on-device twin component. The
UE-anchored Kirchhoff render runs at $\sim\!100\,$ms per pose on a
single CPU core at $T_{\mathrm{vis}} \sim 10^{4}$ visible
triangles. The body-side $\bm{\gamma}$ calibration is a
ridge-least-squares fit over the $K_{99}$-mode SVD basis and runs
in $\sim\!1\,$ms. The pose-gradient finite difference over the
top-six sensitivity-ranked joints requires seven Kirchhoff
renders, $\sim\!700\,$ms total. The downlink pose-suggestion
payload is $K \log_2(M)$ bits per twin tick at typical $K \leq 4$
and $M = 64$, well within the 5G NR downlink control-channel
budget.

\begin{table}[!t]
\caption{Refresh cadences and per-update wall-clock for the
on-device twin loop.}
\label{tab:cadence}
\centering\footnotesize
\begin{tabular}{@{}l c c@{}}
\toprule
Component & Cadence & Per-update wall-clock \\
\midrule
SMPL-X forward map $\thetab \mapsto \Sigma$ & $\sim\!10\,$Hz & $\sim\!5\,$ms \\
UE-anchored Kirchhoff render & $\sim\!10\,$Hz & $\sim\!100\,$ms \\
$\bm{\gamma}$ ridge LS over $K_{99}$ basis & $\sim\!1\,$Hz & $\sim\!1\,$ms \\
Pose gradient (FD on top-6 joints) & $\sim\!1\,$Hz & $\sim\!700\,$ms \\
DL pose-suggestion payload & per-tick & $K \log_2(M) \approx 24\,$bits \\
UL pilot for residual & per-slot & embedded in CSI feedback \\
\bottomrule
\end{tabular}
\end{table}

\subsection{Privacy delta vs. a model-only baseline}

The SMPL-X pose vector $\thetab$ does not leave the device. The BS
sees only the cascaded channel $\hh^{\mathrm{meas}}$ via uplink
pilots and the body-side calibration coefficient $\gammab \in
\Complex^{K}$ for $K \leq M$. The pose suggestion the BS issues is
a direction in pose space, not a target pose. The reading's per-tick
$\Pabs$ scalar is the only body-side telemetry that leaves the
device for cohort use, opt-in.

\subsection{Limitations}

Three limitations of the present work are explicit.
First, the analysis is single-user. Multi-user co-existence, in which
one user's pose moves redirect another user's multipath, is a
parallel research thread not treated here.
Second, the back-face cull suffices for torso, head, and limb-scale
features at FR2; fingertip features and ear edges need strict
ray-mesh visibility, which raises the per-pose render cost by a
factor of $\sim\!5$ and is detailed in the supplementary information.
Third, the sub-6~GHz extension via the layered Fabry--P\'erot
transmission model introduces frequency-selective effects (the
$\lambda/4$ fat-resonance dip in particular) not present in the
mmWave PO regime; the supplementary information treats this case
explicitly.

\subsection{Future work}

The application-twin framing extends naturally to multi-user
co-existence (where one user's pose moves redirect another user's
multipath; flagged as a future paper), to vehicular twins (where the
parametric model is a chassis rather than a body), and to AR/XR
application semantics (where the rate target is itself a function
of the application's frame budget rather than a static MCS cap).


% =====================================================================
\section{Conclusion}\label{sec:conclusion}
% =====================================================================

We have proposed an on-device digital twin of the user's body that
operates as the wireless-native intelligent agent inside the radio
access network's per-slot control loop. The twin produces, from the
same per-triangle Fresnel--physical-optics machinery, a live SAR
reading and a pose-differentiable mmWave channel; uplink-pilot
calibration via a $K$-mode SVD basis closes the loop bidirectionally.
The central new computational primitive is a UE-anchored Kirchhoff
render of the body-mediated channel coefficient that runs in
$\mathcal{O}(M\,T_{\mathrm{vis}})$ per pose and is the natural dual
of forward radiosity at small receiver counts. The empirical
study finds a naturalness-bounded pose move worth $0.2$~to
$10.5\,$dB of cascaded-SINR gain across the $30$ Munich scenes
(mean $3.1\,$dB on LOS, $3.0\,$dB on NLOS), reaching the
$3\,$dB target in $25\,\%$ of LOS and $29\,\%$ of NLOS scenes
and the $10\,$dB rescue target on $7\,\%$ of NLOS scenes;
the gradient mass concentrates on six torso joints out of $22$;
and the peak-local absorbed power density sits at $0.58\,\%$ of
the ICNIRP reference, two orders of magnitude below the regulatory
ceiling. The on-device loop runs
at $\sim\!1$ to $10\,$Hz on a single CPU core. The body twin is
the wireless-native intelligent agent the special-issue scope
calls for, and the live exposure reading it produces opens a
direct measurement path for the prospective RF-EMF cohort
studies the regulatory framing was originally designed to inform.


% =====================================================================
\appendices
% =====================================================================
\section{Per-path Fresnel transmission operator}\label{app:fresnel}

We collect here the formal definition of the Fresnel transmission
operator $\Fn$ that appears throughout
\cref{sec:physics,sec:kirchhoff}; the derivation from Maxwell's
boundary conditions follows the standard route~\cite{BornWolf}.

At the surface point $\bm{r}$ on triangle $t$, path $n$ has
incidence cosine $\mu_n = \hat{\bm{n}}_t \cdot (-\hat{\bm{k}}_n)$.
For a front-facing path ($\mu_n > 0$) the TE and TM unit vectors are
\begin{equation}\label{eq:TE-TM-basis}
  \hat{\bm{e}}_{s,n} = \frac{\hat{\bm{k}}_n \times \hat{\bm{n}}_t}
                            {\|\hat{\bm{k}}_n \times \hat{\bm{n}}_t\|},
  \qquad
  \hat{\bm{e}}_{p,n} = \hat{\bm{e}}_{s,n} \times \hat{\bm{k}}_n\,,
\end{equation}
and the Fresnel amplitude transmission coefficients are
\begin{equation}\label{eq:fresnel-t}
  t_{s,n} = \frac{2\mu_n}{\mu_n + \xi_n}\,, \qquad
  t_{p,n} = \frac{2\,\ntilde\,\mu_n}{\ntilde^2\,\mu_n + \xi_n}\,,
\end{equation}
with $\xi_n \equiv \sqrt{\ntilde^2 - 1 + \mu_n^2}$ and
$\RE(\xi_n) > 0$. For tissue at 28~GHz, the IT'IS database gives
$\ntilde \approx 4.49 - 1.79j$ for skin.

\begin{definition}[Fresnel transmission operator]\label{def:fresnel-op}
The dimensionless rank-2 operator
\begin{equation}\label{eq:Fn-def}
  \Fn(\bm{r}) =
  \bigl(t_{s,n}\,\hat{\bm{e}}_{s,n}\hat{\bm{e}}_{s,n}^{\!\top}
       + t_{p,n}\,\hat{\bm{e}}_{p,n}\hat{\bm{e}}_{p,n}^{\!\top}\bigr)\,H(\mu_n)
\end{equation}
projects $\psin$ onto the local TE and TM directions and scales each
by its Fresnel transmission coefficient; the Heaviside factor
$H(\mu_n)$ vanishes on back-facing paths.
\end{definition}

The transmitted field at depth $z$ inside the tissue is then
\begin{equation}\label{eq:Etrans}
  \EE_n^{\mathrm{trans}}(\bm{r},z)
  = \Fn(\bm{r})\,\psin\,x_{j(n)}\,
    e^{-jk_0\hat{\bm{k}}_n\cdot\bm{r}}\,
    e^{-jk_0\xi_n z}\,,
\end{equation}
where $x_{j(n)}$ is the precoder weight on the source antenna of
path $n$ and $j$ is the engineering imaginary unit. Writing
$k_0\xi_n = \beta_n - j\alpha_n$ with $\alpha_n > 0$ separates the
amplitude decay rate $\alpha_n$ (skin depth $\delta_n = 1/\alpha_n$)
from the in-tissue phase rate $\beta_n$. The per-triangle integrand
of the UE-anchored Kirchhoff render specialises~\eqref{eq:Etrans}
at $z = 0$ and contracts against the UE steering vector
$\bm{g}_{\mathrm{UE}}(\hat{\bm{\eta}})$.

The operator factorisation~\eqref{eq:Fn-def} is the structural
upgrade over per-element scalar Fresnel calibrations (a single
$\beta\approx 0.70$ that lumps Fresnel transmission with depth
decay): it preserves the polarisation rotation that the BS
precoder $\bm{x}$ can exploit at the per-triangle level, and it
isolates the frequency-dependent depth integral into the universal
coupling constant of \cref{prop:approx2-app} below.

% =====================================================================
\section{Approximations 1 and 2 error budget}\label{app:approxs}

The reduction of the coherent absorption double sum to a single
squared norm rests on two well-bounded approximations,
restated and quantified here. Full proofs are in SI~\S{}S1.

\begin{proposition}[Approximation 1: TM polarisation direction]
\label{prop:approx1-app}
The exact transmitted TM field points along the refracted direction
$\hat{\bm{e}}'_{p,n}$ inside tissue, with
$\sin\theta_{t,n} = \sin\theta_n / |\ntilde|$ by Snell's law. The
Fresnel operator $\Fn$ in~\eqref{eq:Fn-def} replaces
$\hat{\bm{e}}'_{p,n}$ by the incident-side $\hat{\bm{e}}_{p,n}$.
Self-terms are unaffected
($\|\hat{\bm{e}}_{p,n}\|^2 = \|\hat{\bm{e}}'_{p,n}\|^2 = 1$).
The error on TM--TM cross-terms is $O(1/|\ntilde|^2)$, which is
$\lesssim 4\,\%$ for skin at 28~GHz.
\end{proposition}

The physical content is that the in-tissue normal field
$E_\perp \propto \sin\theta_{t,n} = O(1/|\ntilde|)$, so its
contribution to $\|\EE^{\mathrm{trans}}\|^2$ is $O(1/|\ntilde|^2)$.
The high tissue refractive index suppresses normal coupling, not
the polarisation direction itself.

\begin{proposition}[Approximation 2: universal depth coupling]
\label{prop:approx2-app}
Define the normalised depth coupling
\begin{equation}\label{eq:Gamma-app}
  \Gamma_{nn'} \equiv
    \frac{2\sqrt{\alpha_n\alpha_{n'}}}
         {\alpha_n + \alpha_{n'} - j(\beta_{n'} - \beta_n)}\,,
\end{equation}
so that $\Gamma_{nn} = 1$ by construction. Setting
$\Gamma_{nn'} \approx 1$ for all pairs reduces the coherent double
sum (SI~\S{}S1, eq.~(S2)) to a single squared norm,
and incurs a non-conservative error $\leq 0.5\,\%$ on every
cross-term for skin at 28~GHz.
\end{proposition}

The mechanism is that $|\ntilde| \approx 4.83$ for skin compresses
all refraction angles to within $\sim\!10^{\circ}$ of normal, so
$\alpha_n$ and $\beta_n$ are nearly angle-independent and
$\Gamma_{nn'}$ stays close to unity. The error is bounded by
sweeping all incidence-angle pairs.

The combined error budget across the millimetre-wave band is
summarised in \cref{tab:approx-acc}.

\begin{table}[!htbp]
\centering
\caption{Combined error budget for Approximations 1 and 2 on
biological skin, from the IT'IS v5.0 database. App.~1 error is
$\max_\theta \|\hat{\bm{e}}'_{p,n} - \hat{\bm{e}}_{p,n}\|$ projected
onto cross-terms; App.~2 error is $\max_{n,n'}|1 - \Gamma_{nn'}|$.
Computed from the analysis of SI~\S{}S1.}
\label{tab:approx-acc}
\begin{tabular}{rrrr}
\toprule
$f$ (GHz) & $|\ntilde|$ & App.~1 (\%) & App.~2 (\%) \\
\midrule
10  & 5.87 & $\lesssim 3$  & $\leq 0.30$ \\
28  & 4.83 & $\lesssim 4$  & $\leq 0.44$ \\
40  & 4.30 & $\lesssim 5$  & $\leq 0.55$ \\
60  & 3.68 & $\lesssim 7$  & $\leq 0.75$ \\
100 & 3.01 & $\lesssim 11$ & $\leq 1.1$  \\
\bottomrule
\end{tabular}
\end{table}

Both errors increase with frequency as $|\ntilde|$ decreases, but
remain comfortably inside the $\sim\!20\,\%$ tissue-dielectric
uncertainty across the FR2 band; at 28~GHz they are an order of
magnitude smaller than the calibration headroom that the body-side
$\gammab$ regression in \cref{sec:closed-loop} absorbs.

% =====================================================================
\section{Identifiability of the $K$-mode SVD calibration}
\label{app:ident}

We give the truncated-SVD argument behind \cref{prop:ident}.

Let $\Lambda_{\mathrm{KH}} \in \Complex^{M \times T_{\mathrm{vis}}}$
be the per-pilot Kirchhoff regressor matrix introduced in
\cref{sec:closed-loop}, with thin SVD
$\Lambda_{\mathrm{KH}} = \bm{U}\bm{\Sigma}\bm{V}^{\!H}$ and singular
values $\sigma_1 \geq \sigma_2 \geq \cdots \geq \sigma_M > 0$. The
ridge least squares with body-side basis
$\bm{V}_K = \bm{V}_{:,1:K}$ solves
\begin{equation}\label{eq:gamma-ridge}
  \hat{\gammab} = \argmin_{\gammab \in \Complex^K}
    \bigl\|\bm{r} - \Lambda_{\mathrm{KH}}\bm{V}_K\gammab\bigr\|^2
  + \lambda \|\gammab\|^2\,,
\end{equation}
where $\bm{r} = \hat{\hh} - \hBS - \hbody \in \Complex^M$ is the
pilot residual after the LOS and nominal-pose body terms are
subtracted, and $\lambda \geq 0$ is a small ridge.

\textbf{Identifiability.} Substituting the SVD and using
$\bm{V}_K^{\!H}\bm{V}_K = \bm{I}_K$,
\begin{equation}\label{eq:Fb-Vk}
  \Lambda_{\mathrm{KH}}\bm{V}_K = \bm{U}_K\bm{\Sigma}_K\,,
\end{equation}
where $\bm{U}_K = \bm{U}_{:,1:K}$ and
$\bm{\Sigma}_K = \mathrm{diag}(\sigma_1,\dots,\sigma_K)$. The
unregularised normal equations give
$\hat{\gammab} = \bm{\Sigma}_K^{-1}\bm{U}_K^{\!H}\bm{r}$, which is
unique whenever $\sigma_K > 0$, i.e.~whenever the truncated basis
sits in the column space of $\Lambda_{\mathrm{KH}}$. If the residual
$\bm{r}$ lies entirely in $\mathrm{span}(\bm{U}_K)$, the recovery is
exact; if not, the orthogonal complement
$(\bm{I}-\bm{U}_K\bm{U}_K^{\!H})\bm{r}$ is by construction outside
the body-side estimable subspace and is correctly absorbed into the
LOS / antenna calibration $\betab$ in
\cref{sec:closed-loop}.

\textbf{Noise amplification.} Under additive complex Gaussian pilot
noise $\bm{n} \sim \mathcal{CN}(\bm{0},\sigma_p^2 \bm{I}_M)$,
the per-mode estimation variance is
\begin{equation}\label{eq:gamma-var}
  \mathrm{Var}(\hat\gamma_k)
    = \frac{\sigma_p^2}{\sigma_k^2 + \lambda}\,,
\end{equation}
so the fit is well-posed in the per-mode SNR
sense exactly when $\sigma_k^2 \,\mathrm{SNR}_{\mathrm{ul}}\gg 1$.
Equivalently, the $K$-th mode is informative when the per-mode SNR
$\sigma_k^2/\sigma_p^2$ exceeds unity. This justifies the operating
choice $K = K_{99}$, defined as the largest $K$ with
$\sum_{k\leq K}\sigma_k^2 \geq 0.99\,\|\bm{\Sigma}\|_F^2$ and
clipped at $\sigma_K^2 / \sigma_p^2 \geq 1$. The $K_{99}$ basis is
adaptive to the BS--user--body geometry: rank-1 in the far-BS regime
where one body-bounce mode dominates, rank-many in the close-BS
regime where the body subtends a wide solid angle and several body
modes carry comparable spectral mass (operator-spectrum sweep is in
the SI).

\textbf{Lower bound on pilot SNR.} The fit is non-trivial when
$\mathrm{SNR}_{\mathrm{ul}} \geq K\sigma_K^{-2}$, i.e.~when the
per-mode noise variance~\eqref{eq:gamma-var} is small enough that
the $K$ leading modes are jointly resolvable. For the served-user
geometries swept in \cref{sec:numerical}, $K_{99} \in \{1, 2, 4\}$
and $\sigma_K^{-2} \leq 1$, so the operating point of
$20\,$dB UL pilot SNR sits comfortably inside the identifiable
region. \hfill$\square$


% =====================================================================
% Bibliography
% =====================================================================
\begin{thebibliography}{99}
\bibitem{MacCartney2017} G.~R. MacCartney \emph{et al.}, ``Rapid
fading due to human blockage in pedestrian crowds at 5G millimeter-wave
frequencies,'' in \emph{Proc. IEEE GLOBECOM}, 2017.
\bibitem{Maccartney2017blockage} G.~R. MacCartney \emph{et al.},
``Path loss models for 5G millimeter wave propagation channels in
urban microcells,'' \emph{IEEE Access}, vol.~6, 2018.
\bibitem{ICNIRP2020} ICNIRP, ``Guidelines for limiting exposure to
electromagnetic fields (100~kHz to 300~GHz),'' \emph{Health Physics},
vol.~118, no.~5, 2020.
\bibitem{SMPLX2019} G.~Pavlakos \emph{et al.}, ``Expressive body
capture: 3D hands, face, and body from a single image,'' in \emph{Proc.
CVPR}, 2019.
\bibitem{Pavlakos2019VPoser} G.~Pavlakos \emph{et al.}, ``VPoser:
variational human pose prior for body inverse kinematics,''
companion software release, MPI-IS, 2019.
Available: \url{https://github.com/nghorbani/human_body_prior}.
\bibitem{AMASS2019} N.~Mahmood \emph{et al.}, ``AMASS: archive of
motion capture as surface shapes,'' in \emph{Proc.\ ICCV}, 2019.
\bibitem{Hoydis2023Sionna} J.~Hoydis \emph{et al.}, ``Sionna~RT:
differentiable ray tracing for radio propagation modelling,'' in
\emph{Proc. IEEE GLOBECOM Workshops}, 2023.
\bibitem{DiRenzo2020} M.~Di~Renzo \emph{et al.}, ``Smart radio
environments empowered by reconfigurable intelligent surfaces,''
\emph{IEEE J. Sel. Areas Commun.}, vol.~38, no.~11, 2020.
\bibitem{Wu2019} Q.~Wu and R.~Zhang, ``Towards smart and
reconfigurable environment: Intelligent reflecting surface aided
wireless network,'' \emph{IEEE Commun. Mag.}, vol.~58, no.~1, 2020.
\bibitem{Madgwick2011} S.~O.~H. Madgwick \emph{et al.}, ``Estimation
of IMU and MARG orientation using a gradient descent algorithm,'' in
\emph{Proc. IEEE ICORR}, 2011.
\bibitem{TS38214} 3GPP, ``NR; Physical layer procedures for data,''
TS~38.214, Table~5.1.3.1-2, 2024.
\bibitem{Hochwald2014} B.~M. Hochwald \emph{et al.}, ``A general
framework for SAR-aware MIMO transmission,'' \emph{IEEE Trans.
Wireless Commun.}, vol.~13, no.~3, 2014.
\bibitem{Ying2015} D.~Ying \emph{et al.}, ``Hermitian
characterization of the maximum rate of energy reception under SAR
constraint,'' \emph{IEEE Trans. Antennas Propag.}, vol.~63, no.~12,
2015.
\bibitem{Ying2017} D.~Ying \emph{et al.}, ``Closed-form approximations
for SAR-constrained beamforming,'' \emph{IEEE Trans. Antennas
Propag.}, vol.~65, no.~7, 2017.
\bibitem{Roosli2010} M.~R\"o\"osli, ``Health effects of mobile phone
use,'' \emph{Bioelectromagnetics}, 2010.
\bibitem{Roosli2021} M.~R\"o\"osli \emph{et al.}, ``Cosmos: Mobile
phone use and brain tumour risk in a prospective cohort study,''
\emph{Environ. Int.}, 2024.
\bibitem{WydaeghePB2024} R.~Wydaeghe \emph{et al.}, ``Pseudo-Brewster
collapse on biological tissue at mmWave,'' under review at
\emph{IEEE Trans. Antennas Propag.}, 2026.
\bibitem{BornWolf} M.~Born and E.~Wolf, \emph{Principles of Optics},
7th ed. Cambridge, U.K.: Cambridge Univ. Press, 1999.
\bibitem{Chen2024DTN} M.~Chen \emph{et al.}, ``Digital twin
networks: a survey,'' \emph{IEEE Internet Things J.}, 2024.
\textbf{[verify at draft-review]}
\bibitem{Saad2025DTN} W.~Saad \emph{et al.}, ``Digital twins for
6G wireless networks,'' \emph{Proc. IEEE}, 2025. \textbf{[verify]}
\bibitem{Hashash2025DTN} O.~Hashash \emph{et al.}, ``Wireless
digital twins for cross-layer optimisation,'' (TBD).
\textbf{[verify]}
\bibitem{CLUEH2025} \emph{CLUE-H consortium: cohort linkage on
mobile-phone RF-EMF exposure and health outcomes,} EU collaborative
project. \textbf{[Robin: insert canonical citation at draft-review.]}
\bibitem{vanWel2024} L.~van~Wel \emph{et al.}, ``Methods for the
RF-EMF exposure side of prospective cohort studies,'' (TBD).
\textbf{[Robin: insert canonical citation at draft-review.]}
\bibitem{general5G} M.~Series, ``IMT vision: framework and overall
objectives of the future development of IMT for 2030 and beyond,''
ITU-R Recommendation M.2160-0, 2023.
\bibitem{HeathLozanoFA2018} R.~W. Heath~Jr.\ and A.~Lozano,
\emph{Foundations of MIMO Communication}. Cambridge University
Press, 2018.
\bibitem{BrusselsArrete2007} ``Arr\^et\'e de la R\'egion de
Bruxelles-Capitale relatif \`a la protection de l'environnement
contre les \'eventuels effets nocifs et nuisances provoqu\'es par
les radiations non ionisantes,'' 2007.
\bibitem{BrusselsArrete2024} ``Arr\^et\'e de la R\'egion de
Bruxelles-Capitale (2024 update),'' 2024.
\bibitem{ItalyDL} ``Italian decree on RF-EMF reference levels for
mobile networks,'' 2024.
\bibitem{DIP2018} Y.~Huang \emph{et al.}, ``Deep inertial poser:
learning to reconstruct human pose from sparse inertial
measurements in real time,'' \emph{ACM Trans.\ Graph.}, 2018.
\bibitem{TransPose2021} X.~Yi \emph{et al.}, ``TransPose: real-time
3D human translation and pose estimation with six inertial
sensors,'' \emph{ACM Trans.\ Graph.}, 2021.
\bibitem{Mollyn2023IMUPoser} V.~Mollyn \emph{et al.}, ``IMUPoser:
full-body pose estimation using IMUs in phones, watches, and
earbuds,'' in \emph{Proc.\ ACM CHI}, 2023.
\bibitem{Xu2024MobilePoser} V.~Xu \emph{et al.}, ``MobilePoser:
real-time full-body pose estimation and 3D human translation
from IMUs in mobile consumer devices,'' in \emph{Proc.\ ACM
UIST}, 2024.
\bibitem{SwissNISV2000} ``Verordnung \"uber den Schutz vor
nichtionisierender Strahlung (NISV) / Ordonnance sur la
protection contre le rayonnement non ionisant (ORNI),''
Swiss Federal Council, SR 814.710, 2000 (and subsequent
amendments).
\bibitem{Christ2010VF} A.~Christ \emph{et al.}, ``The Virtual
Family --- development of surface-based anatomical models of
two adults and two children for dosimetric simulations,''
\emph{Phys.\ Med.\ Biol.}, vol.~55, no.~2, pp.~N23--N38, 2010.
\end{thebibliography}

\end{document}
